\chapter{Simulation 實驗與結果}

本章說明評估資料與標註、對照組與假設、時間回放與延誤注入、評估指標、統計方法與風險。本章所列數字皆為計畫值，實驗尚未執行，目前沒有任何結果數值，結果將於完成後依實際測量填入。


\section{資料與標註}

評估以臺中市為主；納入的餐廳區域與公車路線範圍將於資料檢查後決定。

餐廳資料取自數支 YouTube 美食影片的逐字稿，萃取餐廳名稱、地點、菜色、價位與評價等關鍵字與屬性，並以影片發布日期作為資料時間；另蒐集影片留言，並透過官方 Places API 將餐廳對應到 Google 地圖地點，知識圖譜保存 Place ID 作為對齊鍵；依官方政策~\cite{google_places_policy}，Places 內容不得預先擷取或儲存，本研究因此不使用 Google 評論與地圖內容，Places API 僅用於取得 Place ID。各筆資料保留來源與資料時間。影片的選取準則（頻道、發布期間、涵蓋區域）將於蒐集前明定並公開。營業時間與價格在逐字稿中可能缺漏或過時，其標準答案以店家官方資訊或抽樣的現場、電話核對為準，Google 地圖僅供人工參考、不記錄其內容。納入評估的餐廳以影片提及的臺中市餐廳為主，初步規劃約 100 間，實際數量取決於可萃取者。


查詢以模板與人工撰寫混合產生，依「限制類型 × 用餐時段」分層：限制類型包含預算、營業時間、步行上限、搭乘條件（直達或可轉乘一次）與飲食偏好，用餐時段包含午餐、晚餐與宵夜，讓硬限制在不同時段實際發生違反。初步規劃先導 50 筆、正式 200 筆（其中 60 筆用於調整、140 筆用於測試），樣本數將依先導的違反率估計進行檢定力評估後調整。標註分兩層：硬限制（營業、預算、班次是否存在）以資料來源與即時檢查建立可驗證的標準答案，不依賴評分者主觀判斷；只有偏好相關性由兩位評分者各自獨立標註，並報告評分者一致性。查詢也納入「搭公車去也搭公車回」「晚餐、需要回程末班」等高確定性限制的情境。


\section{對照組與假設}

各研究假設的比較對象與評估指標整理於表~\ref{table:hypotheses}。

\newcolumntype{Y}[1]{>{\raggedright\arraybackslash\hspace{0pt}}p{#1}}

\begin{table}[H]
    \centering
    \caption[研究假設、比較對象與評估指標]{研究假設、比較對象與評估指標}\label{table:hypotheses}
    \renewcommand{\arraystretch}{1.3}
    \setlength{\tabcolsep}{5pt}
    \normalsize
    \begin{tabular}{|Y{\dimexpr 0.1200\textwidth-2\tabcolsep\relax}|Y{\dimexpr 0.3800\textwidth-2\tabcolsep\relax}|Y{\dimexpr 0.5000\textwidth-2\tabcolsep\relax}|}
        \hline
        \textbf{假設} & \textbf{比較對象} & \textbf{看的指標} \\ \hline
        H1 & 關鍵字找餐廳、向量、混合三種召回 & 排名指標（於先導階段確定） \\ \hline
        H2a & 完整證據驗證 vs. AI 直接排名 & 違反限制率、未知誤承諾率；同時報告覆蓋率與拒答率 \\ \hline
        H2b（探索） & 圖譜 vs. 關聯式資料庫版本 & 查詢正確性與證據完整度（指標定義於先導階段確定） \\ \hline
        H3 & 以預估抵達時間檢查營業時間 vs. 以出發時間檢查 & 營業違反率（到店時已打烊的比例），依用餐時段分層報告 \\ \hline
        H4 & 加入即時到站查核 vs. 只用旅運規劃回傳的時間 & 無法上車或轉乘接不上的比例；「不確定」標示的正確率（以現場驗證為標準） \\ \hline
        H5（探索） & 餐廳與搭乘方案合併排序 vs. 先選餐廳再配公車 & 最終方案的總旅程時間、等車時間與違反率；排名因搭乘方案而改變的比例 \\ \hline
        H6 & 行程中監控並重新規劃 vs. 不監控的靜態方案 & 偵測率、誤報率、提醒領先時間與替代方案可搭乘比例（於重播資料上計算，現場結果僅描述） \\ \hline
    \end{tabular}
\end{table}


H2b 與 H5 為探索性比較，不據以宣稱優越；圖譜與關聯式資料庫的比較須以相同事實與規則進行，僅刪除圖譜卻同時刪去事實，無法證明圖譜的因果效益。


\section{連續輪詢與時間回放}

除查詢評估外，另設兩項資料蒐集。其一，連續一至兩週以固定間隔輪詢選定路線的預估到站（N1）端點，統計有即時預估的比例、各狀態碼的占比、缺漏與更新延遲分佈，作為「有即時依據」能否成立的描述性證據，輪詢頻率須符合 TDX 依訂閱方案設定的流量限制。其二，小規模現場驗證，抽取約 30 趟實際搭乘（計畫值），記錄推薦班次是否在預估時間窗內到站、轉乘是否接得上，這也是取得實際到站標準答案的方式。

行程中監控的驗證分三層（表~\ref{table:monitor_validation}），原因是約 30 趟的現場驗證未必遇到延誤。

\newcolumntype{Y}[1]{>{\raggedright\arraybackslash\hspace{0pt}}p{#1}}

\begin{table}[H]
    \centering
    \caption[行程中監控的三層驗證]{行程中監控的三層驗證}\label{table:monitor_validation}
    \renewcommand{\arraystretch}{1.3}
    \setlength{\tabcolsep}{5pt}
    \normalsize
    \begin{tabular}{|Y{\dimexpr 0.1400\textwidth-2\tabcolsep\relax}|Y{\dimexpr 0.4000\textwidth-2\tabcolsep\relax}|Y{\dimexpr 0.4600\textwidth-2\tabcolsep\relax}|}
        \hline
        \textbf{層次} & \textbf{做法} & \textbf{限制與用途} \\ \hline
        重播 & 以連續輪詢所存的 N1 紀錄，找出預估到站時間跳升、StopStatus 轉為 2 至 4 的延誤事件，依紀錄時序重播監控邏輯，與不監控的靜態方案比較偵測率、誤報率、提醒領先時間與替代方案可搭乘比例 & 不使用 A1、A2 動態資料，到站時刻僅能由 N1 紀錄推定，只檢驗監控邏輯相對於 N1 資料的行為，不能證明 N1 預估與實際到站相符 \\ \hline
        延誤注入 & 在重播資料上人為加入不同幅度的延誤與缺漏，以 10 分鐘門檻與 10 分鐘查詢間隔為基準，檢驗門檻與間隔變動下的敏感度 & 結果僅說明方法的行為，不視為真實成效 \\ \hline
        現場 & 優先在尖峰時段等較易延誤的情境抽樣約 30 趟（計畫值），記錄推薦班次是否在預估時間窗內到站、轉乘是否接得上，並記錄被動遇到的延誤次數 & 實際到站的標準答案僅來自現場人工記錄；若遇到的延誤事件不足，結果僅作描述、不做推論檢定，並列為研究限制 \\ \hline
    \end{tabular}
\end{table}


第一層重播的延誤事件，係指預估到站時間跳升或 StopStatus 轉為 2 至 4 者；提醒領先時間為提醒時刻到 N1 紀錄推定之車輛到站時刻的差。所有注入的故障均標示為模擬，輪詢頻率須符合 TDX 依訂閱方案設定的流量限制。


\section{評估指標}

各研究問題對應的指標與分母整理於表~\ref{table:metrics_rq}。

\newcolumntype{Y}[1]{>{\raggedright\arraybackslash\hspace{0pt}}p{#1}}

\begin{table}[H]
    \centering
    \caption[評估指標與研究問題對應]{評估指標與研究問題對應}\label{table:metrics_rq}
    \renewcommand{\arraystretch}{1.3}
    \setlength{\tabcolsep}{5pt}
    \normalsize
    \begin{tabular}{|Y{\dimexpr 0.1300\textwidth-2\tabcolsep\relax}|Y{\dimexpr 0.4000\textwidth-2\tabcolsep\relax}|Y{\dimexpr 0.4700\textwidth-2\tabcolsep\relax}|}
        \hline
        \textbf{研究問題} & \textbf{指標} & \textbf{定義與分母} \\ \hline
        RQ1 & Recall@K、MRR、nDCG（實際排名指標於先導階段確定） & 比較關鍵字、向量、混合三種召回；無相關餐廳的查詢另列，不把其 nDCG 設為完美 \\ \hline
        RQ2 & 限制違反率 & 確有一項違反的「宣稱可行」方案數 ÷ 所有宣稱可行方案數 \\ \hline
        RQ2 & 未知誤承諾率 & 必要證據不足（判為不確定）卻聲稱通過的數量 ÷ 所有宣稱可行方案數 \\ \hline
        RQ2 & 覆蓋率與拒答率 & 至少回傳一筆已驗證方案的查詢 ÷ 所有測試查詢；與違反率同時報告，避免以拒答換取低錯誤 \\ \hline
        RQ3 & 營業時間違反率 & 到店時已打烊的推薦數 ÷ 所有推薦數，依用餐時段分層報告 \\ \hline
        RQ4 & 無法上車率、不確定標示正確率 & 無法上車或轉乘接不上的方案比例；「不確定」標示是否正確，以現場驗證為標準 \\ \hline
        RQ5 & 總旅程時間、等車時間、違反率 & 最終方案的各項時間與違反率；排名因搭乘方案而改變的比例 \\ \hline
        RQ6 & 偵測率、誤報率、提醒領先時間、替代方案可搭乘比例 & 於重播資料上計算，現場結果僅描述（見表~\ref{table:monitor_validation}） \\ \hline
    \end{tabular}
\end{table}


違反率與未知誤承諾率在零輸出時記為不適用，並同時報告覆蓋率與拒答率。


\section{統計與可重現性}

對同一批查詢逐題配對比較，報告 95\% 信賴區間（bootstrap），多重比較以 Holm 校正。負向結果同樣討論，不僅報告成功的比較。同一情境的改寫與模板歸為同一組再切分開發集與測試集，不把同一查詢的多個方案視為獨立樣本。

實驗保存程式碼版本、模型與 embedding 版本、提示與 schema、權重、隨機種子、切分清單與原始回應；正式測試查詢不得用於調整權重或標籤。


\section{風險、限制與降級}

本專題的限制可分為範圍、資料與方法三類，各項風險與因應整理於表~\ref{table:risk_mitigation}。範圍上，不處理訂位、排隊與菜色庫存，只涵蓋公車，不含其他大眾運輸或自行車；所有效能提升都是待驗證的假設，他處研究的結果不能直接證明本場域的成效。

\newcolumntype{Y}[1]{>{\raggedright\arraybackslash\hspace{0pt}}p{#1}}

\begin{table}[H]
    \centering
    \caption[風險與因應]{風險與因應}\label{table:risk_mitigation}
    \renewcommand{\arraystretch}{1.3}
    \setlength{\tabcolsep}{5pt}
    \normalsize
    \begin{tabular}{|Y{\dimexpr 0.2600\textwidth-2\tabcolsep\relax}|Y{\dimexpr 0.7400\textwidth-2\tabcolsep\relax}|}
        \hline
        \textbf{風險} & \textbf{因應} \\ \hline
        餐廳資料偏誤與授權 & 影片只介紹特定餐廳，存在選擇偏誤，逐字稿萃取可能出錯；Places 內容依政策不得儲存，不使用 Google 評論與地圖內容，標準答案改由店家官方資訊與抽樣核對取得，涵蓋範圍較小 \\ \hline
        營業時間與價格缺漏或過期 & 缺漏時判為不確定，不說符合；標準答案以店家官方資訊或抽樣核對為準 \\ \hline
        預估到站不準、缺獨立真值 & 預估到站不保證準確，且不使用 A1、A2 動態資料；重播層僅檢驗監控邏輯，實際到站標準答案僅來自現場記錄 \\ \hline
        預估只涵蓋短時間 & 轉乘後的班次可能超出即時預估範圍，只以班距估計並判為不確定 \\ \hline
        識別碼對不齊 & 站牌、路線、餐廳三類資料以識別碼或名稱對照，路線以 RouteUID 完全比對，對不齊時標為不確定 \\ \hline
        旅運規劃是否使用即時資料未確認 & 旅運規劃結果僅作候選方案，出發時再以 N1 即時預估查核 \\ \hline
        簡化假設與範圍 & 轉乘步行時間、站間行駛時間等推算含簡化假設；只涵蓋臺中市公車，單一地區與人工標註無法支持跨城市結論；正式報告須包含失敗案例與負向結果 \\ \hline
    \end{tabular}
\end{table}

